\subsection{Adaptación de modelos de lenguaje a dominios especializados}
Adaptar un modelo de lenguaje a un dominio especializado suele plantearse, en primer lugar, como un problema de datos: reunir textos y ejemplos que lo acerquen al vocabulario, los temas y las formas de interacción propias del área, una estrategia respaldada por trabajos que encontraron mejoras en tareas específicas al continuar el entrenamiento de modelos generales con datos del dominio \cite{gururangan2020dontstop}. Esta adaptación puede tomar distintas formas. El entrenamiento con textos del área expone al modelo a sus conceptos y expresiones habituales, mientras que el ajuste con pares de consulta y respuesta le muestra cómo se plantean las preguntas y qué tipo de respuesta se espera. En medicina, esta segunda vía resulta relevante para las consultas abiertas, donde el modelo debe responder a situaciones formuladas por personas con necesidades distintas. ChatDoctor exploró ese enfoque mediante el ajuste fino supervisado de un modelo general con conversaciones médicas \cite{li2023chatdoctor}. Sin embargo, contar con ejemplos del dominio no basta para determinar si una respuesta generada es adecuada. En una consulta clínica, el modelo puede emplear términos médicos y mantener el estilo de los datos de entrenamiento, pero omitir información necesaria, perder de vista el contexto de la pregunta o comunicar mal una recomendación \cite{arora2025healthbench}. En este punto se advierte una limitación del ajuste fino supervisado (SFT), ya que el modelo aprende a partir de textos de referencia, pero no recibe por ese mecanismo una señal explícita para distinguir los grados de calidad entre distintas respuestas posibles a una misma consulta \cite{ouyang2022instructgpt}.

